% This section is outside the frozen primary evaluation protocol.
\section{Inference harness and serving: a separate dev-only extension}
\label{sec:harness-serving}

This extension follows checkpoint training. It changes neither the SFT/GRPO objectives nor
the frozen primary evaluation in \cref{sec:eval}. All effectiveness and speed results below
are pending. Its purpose is to distinguish three possible sources of a change in behavior:
learned weights, procedural context at inference, and decoding implementation.

\subsection{Versioned rules and progressive procedural context}
\label{sec:rules-skills}
\Cref{fig:rules-skills-mechanism} shows how the harness supplies procedural context and audits its use. The harness contract extracts existing system/tool instructions and verifier/reward constraints
into \texttt{rules.md}. Versioned rule identifiers link statements to
prompt rendering or enforcement checks. A byte-level test requires the default file-based rendering
to reproduce the existing system prompt exactly; inventory checks bind verifier/reward
identifiers to their implementations without changing outcomes or reward values. Code-only and
prompt-only constraints are recorded as discrepancies, not silently reconciled.


Five versioned procedural skills cover box verification, caries versus deep caries, periapical
lesions, impacted teeth, and a healthy-scan negative survey. Sources accompany each skill;
content is drawn from existing procedures, class definitions, and cited public material,
without test images or test-derived examples. Skills are contextual instructions, not new
learned parameters. With skills disabled, the existing agent is unchanged. In \texttt{index\_only} mode, the prompt lists names and descriptions, and a
\texttt{read\_skill} harness action loads a body on demand. In \texttt{full\_inject} mode, all
bodies are present initially. The harness action is dispatched outside the unchanged eight-tool
diagnostic registry; it does not consume a diagnostic-call slot, but the episode turn bound
still applies. Loaded skill versions, content digests, and turn/timing events are recorded in
optional trace metadata. Old traces remain readable without that metadata.


\begin{figure}[tbp]
\centering
\includegraphics[width=\linewidth]{fig_rules_skills_mechanism}
\caption{Rules and skills in the policy context. (a) Contract entries specify structured
output, FDI tooth identity and complete findings. (b--c) Rule rendering and the five-file skill index
construct the initial context. (d) A policy request is validated and dispatched; the returned skill
body and references are appended to conversation history. (e) Disclosure events record which
content was supplied and when. The excerpt is from the existing caries skill.}
\label{fig:rules-skills-mechanism}
\end{figure}

\subsection{Paired evaluation without reopening the primary protocol}
The extension (\cref{tab:harness-dev}) uses only the fixed, training-derived dev split. A frozen checkpoint is evaluated
under baseline, rules only, rules plus index/on-demand skills, and rules plus full skills.
The same image identifiers, seeds, geometry, dispatcher, tool budgets, and scoring implementation
are used across conditions; full skills are a disclosure control, not an additional training arm.
Rules-only and baseline must render identical prompts; any deviation invalidates the control.
We report pooled exact F1 and the finding-level false-discovery fraction
\[
 F_1^{\mathrm{exact}}=\frac{2TP}{2TP+FP+FN},\qquad
 I_{\mathrm{finding}}=\frac{FP}{TP+FP}.
\]
The latter is the harness JSON field \texttt{invented\_findings\_rate}; it is \emph{not} the
healthy-image invention rate of \cref{sec:eval}. A zero denominator is undefined, not zero.
Paired uncertainty resamples whole shared images. We also report format failures and disclosure events alongside the primary endpoints.
Because these conditions were added after the primary protocol was frozen, they are exploratory
and do not amend that protocol. Dev-split selection cannot establish independent generalization.

\begin{table}[tbp]\centering\small
\caption{Exploratory, paired dev-only harness comparison on one frozen checkpoint. Finding invention is $FP/(TP+FP)$, distinct from the primary healthy-image metric. Pending cells are not measurements.}
\label{tab:harness-dev}
\begin{tabular}{@{}lrr@{}}
\toprule
Context condition & Exact F1 & Finding invention \\
\midrule
Baseline & \todo{pending} & \todo{pending} \\
Rules only (identity control) & \todo{pending} & \todo{pending} \\
Rules + on-demand skills & \todo{pending} & \todo{pending} \\
Rules + full skills & \todo{pending} & \todo{pending} \\
\bottomrule
\end{tabular}
\end{table}


\subsection{Speculative decoding: accuracy and end-to-end efficiency}
\label{sec:specdecode}
Speculative decoding accelerates autoregressive generation by proposing several next tokens
and checking them together with the target model \citep{leviathan2023fast,chen2023accelerating}.
Accepted proposals extend the output without the same number of sequential target-decoding
steps. Its role in this system is confined to the local serving backend
(\cref{fig:speculative-mechanism}): the agent still consumes the image and conversation,
produces a structured turn, invokes the shared dispatcher, and receives real tool observations.
Image prefill and diagnostic-tool execution remain separate costs; proposal overhead and
rejected tokens can offset decoding gains.

The comparison uses one merged, trained checkpoint with speculative decoding disabled or
enabled. Both conditions use the same fixed dev images, prompt construction, greedy decoding,
seeds, tool implementation and budgets. The serving configuration, hardware, software versions
and checkpoint identity accompany the result artifact. When the selected serving configuration
uses the native MTP head, the checkpoint audit checks its tensors against the original base
before execution.

\begin{figure}[tbp]
\centering
\includegraphics[width=\linewidth]{fig_speculative_mechanism}
\caption{Speculative decoding and planned draft training. (a) A trained draft proposes a token
block from committed context. (b) The visibility mask exposes the cached prefix and only earlier
draft positions to each next-token prediction. (c) The symbolic greedy example accepts A and B,
replaces the first mismatch with target token X, and discards the dependent suffix. (d) Planned
draft training uses training contexts and frozen-target continuations, then deploys the draft
alongside the unchanged target. No measured hit rate or speedup is implied.}
\label{fig:speculative-mechanism}
\end{figure}

The planned draft-training stage uses with-tools training contexts and continuations generated
by the trained target. The target remains frozen while draft parameters learn to predict its
continuations. Dev and test images do not enter this stage. The resulting draft supplies proposals
in \cref{fig:speculative-mechanism}; the target continues to verify them. Acceptance and
end-to-end efficiency require measurements after training; no completed draft-training result
is claimed here.

\Cref{tab:serving-efficiency} reports diagnostic accuracy as pooled exact F1, end-to-end
speedup as the baseline's total episode wall time divided by the speculative condition's total
wall time, and proposal hit rate as accepted proposed tokens divided by all proposed tokens.
The hit rate is a decoding statistic, not diagnostic accuracy; it is undefined for the baseline.
Timing includes the complete agent episode so that a decoding improvement is not mistaken for
the same proportional reduction in image-processing or tool time. Separately exposed prefill,
decode and tool counters remain in the result artifact.

Decoded turns, tool actions and final answers are compared on each shared image, together
with pooled exact F1. Divergences remain visible and disqualify an equivalence-qualified
speedup; an aggregate score match alone is insufficient. This finite dev comparison does not
establish token-ID equivalence or equivalence under stochastic sampling. The manuscript table
contains only the explicitly selected speculative configuration, recorded with its provenance,
and the paired baseline; it does not select the fastest setting from the evaluation results.
All unmeasured entries remain pending.

\begin{table}[tbp]\centering\small
\setlength{\tabcolsep}{4pt}
\caption{Accuracy and efficiency of speculative decoding on paired dev episodes. Diagnostic accuracy is pooled exact F1. Speedup is baseline total episode wall time divided by speculative total wall time. Proposal hit rate is accepted proposed tokens divided by all proposed tokens, not diagnostic accuracy. The baseline has no proposals. Failed decoded-output or F1 agreement withholds speedup; missing measurements remain pending.}
\label{tab:serving-efficiency}
\begin{tabular}{@{}lrrrr@{}}
\toprule
Decoding & Exact F1 & Speedup & Hit rate & Agreement \\
\midrule
Baseline & \todo{pending} & \todo{pending} & -- & pending \\
Speculative & \todo{pending} & \todo{pending} & \todo{pending} & pending \\
\bottomrule
\end{tabular}
\end{table}

% No measured plot supplied. Original diagrams remain in place.

